
\subsection{Proposal}

We take a graph, upon which there are players who which to identify if the
graph possesses some property. We wish to use this model to find the minimum
communication required for each player to identify if this property exists.
Each Player will receive some portion of the graph as their input--this can be
either vertices, edges, or both. We mostly study reachability of vertices, in
which the only paths considered valid are paths which use only the portions of
the graph given to the players as inputs.

\subsection{Project Goals}

Therefore, the goal of our project is to consider variations of the above
problem statement, and examine the communication complexity for each. Some
variations we can consider might include the followig

\begin{itemize}
    \item Each player receives as input a set of vertices. An Edge $(u,v)$ is
	    considered traversable if $u$ and $b$ are both in the same players
	    input. What is the communication complexity for identifying
	    reachability of two arbitrary vertices?
    \item Consider the same problem as above, but with more than 2 player.
	    Assume also that the inputs form connected graphs. How does the
	    communication scale with the number of players?
    \item Consider again the same problem as above, but with more than 2
	    players and no constraints on the connectedness of the graphs that
	    their inputs produce. Now how does the communication scale with the
	    number of players?
    \item The game changes sequentially. Players start with perfect knowledge
	    of the problem, but then one or both players gains or loses a
	    vertex. 
\end{itemize}

\subsection{Actionables}

Therefore, this project can be summarized into three specific, actionables:

\begin{enumerate}
    \item Define new models or problems in the field.
    \item Write a paper that outlines the formal, technical communication complexity bounds for a given multi-agent traversal problem with certain constraints.
    \item Give a poster presentation of the key results of the paper, with the intention of making the results applicable to a wider audience.
\end{enumerate}

Therefore, our work will remain theoretically grounded while remaining applicable for the wider computing community.

\subsection{Deliverables}

We have three main deliverables:

\begin{enumerate}
    \item A paper presenting actionable item 1 and 2.
    \item A research poster of our key results, for actionable 3.
    \item A website containing the paper and poster information, for actionable 3.
\end{enumerate}

To be more specific, the paper will take the form of a research paper, where we define the problems we are working on, prove bounds on these problems, and show how they model situations in the field of multi-agent graph traversal. The paper is our primary project output.


We allocate two weeks to reading to ensure we understand the subdomain and are not needlessly duplicating results.

\begin{table}[h!]
    \centering
    \renewcommand{\arraystretch}{1.3}
    \begin{tabularx}{\textwidth}{|c|X|X|X|X|X|}
    \hline
    \textbf{Week} & \textbf{Focus} & \textbf{Ciro} & \textbf{Darren} & \textbf{Ivy} & \textbf{Ryan} \\
    \hline
    
    1 & Field Background Knowledge 
      & Read and present 1) 
      & Read and present 1) 
      & Read and present 2) 
      & Read and present 2) \\
    \hline
    
    2 & Field Background Knowledge
      & Read and present 3) 
      & Read and present 4) 
      & Read and present 3) 
      & Read and present 4) \\
    \hline
    
    3 & Define the simplest case and prove it 
      & Each proves a lemma (incomplete information)
      & Each proves a lemma (randomness)
      & Each proves a lemma (intersecting paths)
      & Each proves a lemma (duplicate endpoints) \\
    \hline
    
    4 & Consider obstacles
      & Each proves a lemma
      & Each proves a lemma
      & Each proves a lemma
      & Each proves a lemma \\
    \hline
    
    5 & Consider weights on edges
      & Each proves a lemma
      & Each proves a lemma
      & Each proves a lemma
      & Each proves a lemma \\
    \hline
    
    6 & Consider more players
      & Each proves a lemma
      & Each proves a lemma
      & Each proves a lemma
      & Each proves a lemma \\
    \hline
    
    7 & Wrap up
      & Writer
      & Figures
      & Equation
      & Figures \\
    \hline
    
    \end{tabularx}
    \end{table}

The annotation 1) refers to \cite{One}, 2) to \cite{Two}, 3) to \cite{Three}, and 4) to \cite{Four}
    
